\documentclass[10pt]{article}

\usepackage[a4paper,margin=1.6cm]{geometry}
\usepackage{amsmath}
\usepackage{booktabs}
\usepackage{enumitem}
\usepackage{microtype}
\usepackage{multicol}
\usepackage[compact]{titlesec}
\usepackage{xcolor}
\usepackage[colorlinks=true,urlcolor=blue!45!black]{hyperref}

\titlespacing*{\section}{0pt}{0.65em}{0.25em}
\titlespacing*{\subsection}{0pt}{0.45em}{0.2em}
\setlist{nosep,leftmargin=*}
\setlength{\parskip}{0.18em}
\setlength{\parindent}{0pt}
\emergencystretch=1.5em

\title{\vspace{-1.6em}\Large\textbf{Soft Identity Constraints for LLM-Written Knowledge Graphs}\\[0.15em]
\normalsize A vision for reliable graph memory under incremental ingestion\vspace{-0.5em}}
\author{\normalsize Vishisht Choudhary \,\textperiodcentered\, \texttt{vishi@toffee.at}}
\date{\normalsize Research vision, July 29, 2026\vspace{-1.3em}}

\begin{document}
\maketitle

\section{Situation}

Knowledge graphs are becoming a memory substrate for AI agents. Systems such as
Graphiti, HippoRAG, GraphRAG, and Mem0 convert conversations and documents into
entities and relations, then retrieve context through graph search and random walks.
Unlike a conventional database, however, these graphs are written by a probabilistic
model operating on partial context. The engine still assumes stable entity identity,
valid edges, consistent schemas, and ordered updates, although no human curator now
guarantees those invariants \cite{zep,hipporag,hipporag2,mem0,graphbench}.

This mismatch produces structural, temporal, and semantic faults: one entity is split
across aliases; distinct entities are merged into a generic hub; relations drift in
name and granularity; stale facts remain active; and provenance is flattened. Evidence
from research and production systems points especially to entity resolution. Reported
duplication reaches 27--36\% in one document study, resolution degrades as context
grows, and production pipelines combine exact matching, embedding retrieval, and LLM
adjudication without claiming the problem is solved
\cite{linkkg,itext2kg,kggen,cai,degrag}. Incremental ingestion is the critical setting for agent memory:
each episode arrives without the global view a human curator would use to enforce
\texttt{UNIQUE}.

\section{Vision: Integrity Mechanisms for Probabilistic Writers}

The long-term goal is an ``immune system'' for machine-written graphs: integrity
mechanisms enforced by the storage and retrieval engine rather than delegated to an
extraction prompt. Each violated invariant should have a measurable fault, a
write-time guard, a reversible repair, and an end-to-end benchmark. Identity is the
first work package because both under-resolution and over-resolution alter the paths
through which an agent assembles evidence.

The central design principle is that uncertain identity should not be forced into a
binary merge decision. On insertion, the engine should retain source nodes and
provenance, retrieve possible aliases, and write provisional \texttt{same\_as} edges
with calibrated confidence. At read time, retrieval may traverse those edges in
proportion to confidence. A correct provisional link reconnects scattered facts; an
incorrect link has bounded influence and can be revoked. This is safer than an
irreversible hard merge and more useful than leaving likely aliases disconnected.
Existing systems already use multi-stage resolution \cite{zep,linkkg,degrag}. This
motivates a separate question: should uncertainty be exposed to retrieval rather than
collapsed into a single merge decision?

For a node $v$ with identity confidence $c_v$, a simple retrieval operator reserves a
bounded share of its outgoing probability for identity transitions:
\[
P(v,\cdot)=(1-\lambda c_v)P_{\mathrm{graph}}(v,\cdot)
            +\lambda c_v P_{\mathrm{identity}}(v,\cdot),
\]
where $\lambda$ caps the effect of identity uncertainty. The broader research vision
extends the same pattern to temporal validity, schema equivalence, provenance, and
concurrent writers.

\section{Preliminary Evidence}

Our strongest preliminary result shows one identity failure and a partial repair. We
used MuSiQue, a public collection of questions whose answers require combining facts
from two or more source paragraphs, meaning short sections of the original documents
\cite{musique}. The dataset identifies which
paragraphs contain the required facts. Claude and Codex independently converted the
same 1,048 paragraphs into two graphs. We evaluated macro-averaged supporting-passage
Recall@5: the fraction of gold supporting paragraphs present in the top five retrieved
results, averaged over questions.

We restricted the test to questions with two required paragraphs that shared an
entity. In the intact graph, one entity node was connected to facts from both
paragraphs. We then simulated an identity error by splitting it into two disconnected
nodes: the first kept only the edges extracted from paragraph A, and the second kept
only those from paragraph B. We began retrieval at the first node and measured whether
paragraph B appeared in the top five results. Finally, we connected the two nodes with
a known-correct \texttt{same\_as} edge and repeated retrieval. The table reports the
score before the split, after the split, and after reconnection. Its final column is
the fraction of the lost score restored by the link:

\begin{center}
\small
\begin{tabular}{lrrrr}
\toprule
Extractor & Intact graph & Split graph & With link & Loss recovered \\
\midrule
Claude Opus 5 & 75.0\% & 15.0\% & 64.2\% & 81.9\% \\
Codex GPT-5.6-sol & 77.3\% & 10.0\% & 71.8\% & 91.9\% \\
\bottomrule
\end{tabular}
\end{center}

\newpage
The changes were consistent across questions: repeatedly resampling them produced 95\%
uncertainty ranges that always showed a loss after splitting and a gain after linking.
This demonstrates the mechanism under controlled conditions, not yet a common
production failure. We created the split and supplied the correct link ourselves. We
have not measured how often real ingestion creates duplicate nodes or whether an
automatic resolver finds the right link. The focused test was also designed after a
broader experiment found no effect, so it needs an independent replication fixed in
advance.

We separately tested ``supernodes,'' nodes with unusually many edges. Some graph
retrievers rank nodes by repeatedly simulating a walk along edges \cite{hipporag,hipporag2}.
A supernode may then appear important for many unrelated questions. In one 40-part
conversation, the same highly connected people ranked near the top for many different
queries. Penalizing nodes with many and varied edges reduced this concentration, but
did not reliably improve retrieval of the correct source paragraphs. The supposed bad
supernodes were often legitimate people or places, and a simple degree penalty worked
as well as the proposed relation-aware score. Related work also finds that extraction
errors do not change graph statistics uniformly, and that graph retrieval benefits
depend strongly on the task and construction pipeline \cite{cai,graphbench}.

This negative result may reflect scale: the study used one conversation and 60
multi-source questions. A larger and more diverse experiment could reveal a smaller
effect. Scale is not assumed to be the explanation, however. The proposed score may
be wrong, or direct text similarity may bypass the graph error. A larger study with
its analysis fixed in advance should distinguish these possibilities.

\section{Research Programme}

\textbf{Aim 1: establish the natural fault.} Compare the same extractor and corpus
under (A) batch extraction with global context, (B) blind incremental extraction,
and (C) incremental extraction with resolution. Repeat extraction to measure sampling
instability. Manually adjudicate mention-to-entity clusters on a fixed sample. Primary
outcomes are the duplicate-node rate, how an entity's facts are divided among its
duplicates, and the percentage of required sources found in the first five results.
Questions answerable from one source are the control. The prediction is conditional:
fragmentation should hurt when it
breaks a necessary evidence path, not uniformly across all questions
\cite{linkkg,itext2kg,degrag}.

\textbf{Aim 2: build reversible resolution.} Implement a cost-aware cascade:
normalized exact matching, semantic similarity to propose candidates, contextual
scoring, and LLM adjudication only for ambiguous cases. Preserve original nodes and attach calibrated
confidence to provisional identity edges. Compare no resolution, hard merge, soft
identity flow, and known-correct links under identical extraction and retrieval budgets.

\textbf{Aim 3: test safety and generality.} Inject incorrect identity links to measure
false bridging and compare graceful degradation with destructive hard merging. Run on
LoCoMo and MuSiQue, then confirm on an untouched multi-document or long-memory
benchmark. Replicate across at least two extractor families. Tune on development data
and evaluate once on questions set aside before tuning; report uncertainty across
questions.

\textbf{Success gates.} Natural incremental ingestion must produce more fragmentation
than batch ingestion, with retrieval loss concentrated on bridge-dependent questions.
The practical resolver should recover at least half of the improvement achieved by
known-correct links, create fewer false bridges than hard merging, and average fewer than one LLM
adjudication call per episode. If natural harm is not observed, the result will be
reported as a robustness mechanism rather than evidence of a common production fault.

\section{Expected Contribution and Scope}

The project contributes: (1) a controlled characterization of identity failure under
incremental graph construction; (2) an evaluation focused on questions that require
an identity connection; and (3) a reversible, confidence-weighted retrieval rule for
identity links. The immediate deliverable is deliberately narrow: prove the
natural fault, then test whether soft resolution repairs it safely. Temporal
consistency, schema drift, provenance policy, and multi-agent concurrency remain
future work within the larger integrity-mechanism vision.

\begin{multicols}{2}
\begin{thebibliography}{11}
\tiny
\raggedright
\setlength{\itemsep}{0pt}
\bibitem{hipporag} Guti\'errez et al. \emph{HippoRAG: Neurobiologically Inspired Long-Term Memory for LLMs}. \href{https://arxiv.org/abs/2405.14831}{arXiv:2405.14831}, 2024.
\bibitem{hipporag2} Guti\'errez et al. \emph{From RAG to Memory: Non-Parametric Continual Learning for LLMs}. \href{https://arxiv.org/abs/2502.14802}{arXiv:2502.14802}, 2025.
\bibitem{zep} Rasmussen et al. \emph{Zep: A Temporal Knowledge Graph Architecture for Agent Memory}. \href{https://arxiv.org/abs/2501.13956}{arXiv:2501.13956}, 2025.
\bibitem{mem0} Chhikara et al. \emph{Mem0: Building Production-Ready AI Agents with Scalable Long-Term Memory}. \href{https://arxiv.org/abs/2504.19413}{arXiv:2504.19413}, 2025.
\bibitem{linkkg} Meher et al. \emph{LINK-KG: LLM-Driven Coreference-Resolved Knowledge Graphs}. \href{https://arxiv.org/abs/2510.26486}{arXiv:2510.26486}, 2025.
\bibitem{itext2kg} Lairgi et al. \emph{iText2KG: Incremental Knowledge Graphs Construction Using LLMs}. \href{https://arxiv.org/abs/2409.03284}{arXiv:2409.03284}, 2024.
\bibitem{kggen} Mo et al. \emph{KGGen: Extracting Knowledge Graphs from Plain Text with Language Models}. \href{https://arxiv.org/abs/2502.09956}{arXiv:2502.09956}, 2025.
\bibitem{cai} Cai and O'Connor. \emph{Understanding the Effect of KG Extraction Error on Downstream Graph Analyses}. \href{https://arxiv.org/abs/2506.12367}{arXiv:2506.12367}, 2025.
\bibitem{degrag} Zheng et al. \emph{Less is More: Denoising Knowledge Graphs for Retrieval Augmented Generation}. \href{https://arxiv.org/abs/2510.14271}{arXiv:2510.14271}, 2025.
\bibitem{graphbench} Xiang et al. \emph{When to Use Graphs in RAG: A Comprehensive Analysis}. \href{https://arxiv.org/abs/2506.05690}{arXiv:2506.05690}, 2025.
\bibitem{musique} Trivedi et al. \emph{MuSiQue: Multihop Questions via Single-Hop Question Composition}. \href{https://arxiv.org/abs/2108.00573}{arXiv:2108.00573}, 2021.
\end{thebibliography}
\end{multicols}

\end{document}
